\section{Boîte à outils: quatre définitions et les arches}\label{sec:toolbox}

Cette section contient toute la théorie dont la suite a besoin: quatre définitions, et les cinq lemmes qui forment le registre des arches.

\subsection{Formes et préfaisceaux}

Une petite catégorie $C$ est une \emph{forme}. Un \emph{préfaisceau} sur $C$ est un foncteur $C^{\op}\to\Set$, en Lean \texttt{Cᵒᵖ ⥤ Type u}: un ensemble $F(U)$ pour chaque objet et, pour chaque flèche $f : V\to U$, une application de \emph{restriction} $F(U)\to F(V)$, $x\mapsto\res xf$, avec $\res x{\id} = x$ et $\res{(\res xg)}f = \res x{g\circ f}$. Les préfaisceaux sur $C$ forment le topos $\PSh(C)$. Voici les exemples que nous utilisons:

\begin{center}\small
\begin{tabularx}{\linewidth}{>{\raggedright\arraybackslash}p{4.3cm}>{\raggedright\arraybackslash}X}
\toprule
forme $C$ & préfaisceaux sur $C$\\
\midrule
$\B M$: un objet, flèches $=$ éléments d'un monoïde $M$, $g\circ f = gf$ & $M$-ensembles à droite: un ensemble $X$ muni de $x\cdot m$, avec $x\cdot 1 = x$, $(x\cdot g)\cdot f = x\cdot gf$; dans Mathlib, \texttt{MulAction Mᵐᵒᵖ X}\\
$\B(\Sigma^*)$, $\Sigma^*$ les mots sur un alphabet & automates déterministes: états $X$, transitions $x\cdot a$ (sans état initial ni états finaux)\\
$\B(\mathbb N,+)$ & systèmes dynamiques discrets: un ensemble muni d'une application $\tau : X\to X$\\
la catégorie libre d'un carquois & représentations du carquois opposé: un ensemble par sommet, une application par flèche\\
$V\rightrightarrows E$ & graphes orientés\\
$\Delta$, $\Delta_+$ & ensembles simpliciaux, ensembles simpliciaux augmentés\\
$\omega = (\mathbb N,\le)$ & suites $X_0\leftarrow X_1\leftarrow\cdots$ (le \q{topos des arbres} de la récursion gardée)\\
\bottomrule
\end{tabularx}
\end{center}

Pour un monoïde, la restriction de $x$ le long de $m$ est $x\cdot m$.

\subsection{Pièges et négation}

Un \emph{sous-préfaisceau} $A\subseteq F$, que nous appelons aussi un \emph{piège}, est une famille de sous-ensembles stable par restriction. Pour un $M$-ensemble, c'est une partie stable par l'action. Les sous-préfaisceaux sont les sous-objets du topos et jouent le rôle de propositions. Leur négation n'est pas le complémentaire ensembliste (qui n'est en général pas un piège), mais le plus grand sous-préfaisceau disjoint de $A$:

\begin{lemma}\label{lem:neg}
On a $x\in\neg A$ si et seulement si aucune restriction de $x$ n'appartient à $A$. Par conséquent, $x\in\neg\neg A$ si et seulement si toute restriction de $x$ admet une restriction ultérieure dans $A$.
\end{lemma}

\begin{proof}
L'ensemble de ces $x$ est stable par restriction, disjoint de $A$, et contient tout piège disjoint de $A$. La seconde phrase est la première appliquée à $\neg A$.
\end{proof}

\begin{definition}
Le topos $\PSh(C)$ est \emph{booléen} si $A\cup\neg A = F$ pour tout $A\subseteq F$ (tiers exclu), et \emph{de De Morgan} si $\neg A\cup\neg\neg A = F$ pour tout $A\subseteq F$ (tiers exclu faible). D'après le lemme~\ref{lem:neg}, De Morgan signifie: \emph{pour tout préfaisceau $F$, tout $A\subseteq F$ et tout élément $x$, ou bien aucune restriction de $x$ n'appartient à $A$, ou bien toute restriction de $x$ admet une restriction ultérieure dans $A$.}
\end{definition}

Informellement, De Morgan est une propriété de \emph{stabilité}: \q{$A$ est accessible} ne peut pas se perdre en restreignant davantage. Être booléen est plus fort: \q{$x\in A$} est décidé.

\subsection{La condition d'Ore}

La catégorie $C$ vérifie la \emph{condition d'Ore} si toute paire de flèches de même but $f : Y\to X$, $g : Z\to X$ se complète en un carré commutatif $f\circ h = g\circ k$. Pour $\B M$: pour tous $a,b\in M$, il existe $c,d$ avec $ac = bd$, c'est-à-dire $aM\cap bM\neq\emptyset$ (deux éléments quelconques ont un multiple à droite commun). Elle est vraie pour les groupes, les monoïdes commutatifs, les monoïdes à zéro, les catégories ayant un objet initial ou des produits fibrés, et les chaînes.

\subsection{Les cinq arches}

La bibliothèque \lp{} démontre en Lean~4 les cinq ponts suivants. Nous citons les déclarations principales; les deux premières sont données avec leur énoncé exact.

\begin{bridge}[Booléen]\label{br:bool}
$\PSh(C)$ est booléen si et seulement si $C$ est un groupoïde.
\begin{leancode}
theorem isBooleanTopos_presheaf_iff_isGroupoid (C : Type u) [SmallCategory C] :
    IsBooleanTopos (Cᵒᵖ ⥤ Type u) ↔ IsGroupoid C
\end{leancode}
\end{bridge}

\begin{bridge}[De Morgan]\label{br:dm}
$\PSh(C)$ est de De Morgan si et seulement si $C$ vérifie la condition d'Ore.
\begin{leancode}
theorem isDeMorganTopos_presheaf_iff_rightOreCondition (C : Type u)
    [SmallCategory C] :
    IsDeMorganTopos (Cᵒᵖ ⥤ Type u) ↔ GrothendieckTopology.RightOreCondition C
\end{leancode}
Sa preuve est courte. Si $\res xf\in A$ et si $g$ est une autre flèche de but $X$, on complète le carré $f\circ h = g\circ k$; alors $\res{(\res xg)}k = \res{(\res xf)}h\in A$. Réciproquement, on teste sur le préfaisceau $\y X = \Hom(-,X)$, le piège engendré par $f$, et l'élément $\id_X$.
\end{bridge}

\begin{bridge}[Morita]\label{br:morita}
On a $\PSh(C)\simeq\PSh(D)$ si et seulement si les \emph{enveloppes de Karoubi} $\Kar(C)$ et $\Kar(D)$ sont équivalentes. $\Kar(C)$ ajoute un objet pour chaque flèche idempotente $e = e\circ e$; pour $\B M$, ses objets sont les idempotents de $M$, et $e\cong f$ dans $\Kar(\B M)$ si et seulement s'il existe $a\in fMe$, $b\in eMf$ avec $ba = e$ et $ab = f$. Lean: \lean{moritaEquivalentCat_iff_karoubi}; corollaires \lean{isGroupoid_iff_of_moritaEquivalentCat}, \lean{rightOreCondition_iff_of_moritaEquivalentCat}.
\end{bridge}

\begin{bridge}[Points essentiels]\label{br:pts}
Les foncteurs $\PSh(C)\to\Set$ qui préservent toutes les limites et toutes les colimites (les \emph{points essentiels}, ou \q{foncteurs fibres}) forment une catégorie équivalente à $\Kar(C)^{\op}$. Pour $\B M$, ce sont les foncteurs $X\mapsto Xe = \{x : x\cdot e = x\}$, pour $e$ idempotent, deux d'entre eux étant isomorphes si et seulement si $e\cong f$ dans $\Kar(\B M)$. Lean: \lean{essentialPointEquivKaroubiOp}, \lean{isEssentialPoint_iff}.
\end{bridge}

\begin{bridge}[Sous-topos]\label{br:sub}
Les sous-topos de $\PSh(C)$ (sous-catégories pleines stables par limites dont l'inclusion admet un adjoint à gauche qui préserve les limites finies) correspondent, en renversant l'ordre, aux \emph{topologies de Grothendieck} sur $C$: des choix de cribles \q{couvrants} sur chaque objet, stables par changement de base et transitifs. Lean: \lean{subtoposOrderIso}, \lean{exists_orderIso_ici_of_sheafEquiv}.
\end{bridge}

C'est tout. Chaque énoncé qui suit est l'un de ces cinq lemmes, spécialisé à une forme et déplié à l'aide du lemme~\ref{lem:neg}.

\subsection{Trois arches de plus}\label{sec:morearches}

Trois autres arches ont leur place dans le registre. Elles n'ont été jointes pour l'instant qu'aux tabliers des notes de Grothendieck, ci-dessous; aucun résultat des sections suivantes n'en provient.

\begin{bridge}[Deux valeurs]\label{br:twoval}
Les sous-objets de l'objet final de $\PSh(C)$ (ses \q{ouverts}) correspondent aux cribles de $C$ \cite[IV~8.4.4]{sga4}. Ainsi $\PSh(C)$ n'a pas d'autre sous-objet de l'objet final que $0$ et $1$ si et seulement si, pour deux objets quelconques $a$, $b$ de $C$, il existe une flèche $a\to b$. Lean: \lean{forall_subfunctor_terminal_eq_bot_or_eq_top_iff}.
\end{bridge}

\begin{bridge}[Connexe]\label{br:conn}
$\PSh(C)$ est connexe non vide au sens de \cite[IV~8.7]{sga4} (toute partition de l'objet final a exactement une partie non vide) si et seulement si $C$ est connexe. En termes de sous-objets de l'objet final: $C$ est non vide et les seuls sous-objets complémentés de l'objet final sont $0$ et $1$. Lean: \lean{isConnected_iff_forall_isCompl}.
\end{bridge}

\begin{bridge}[Spatial]\label{br:spatial}
$\PSh(C)$ est équivalent au topos de préfaisceaux d'un préordre, donc au topos des faisceaux sur son espace d'Alexandrov, si et seulement si $C$ est mince. Grothendieck l'énonce comme exercice \cite[p.~26]{grothendieck-161-3}; voir le \S\ref{sec:equivalences}. Lean: \lean{exists_preorder_moritaEquivalentCat_iff_isThin}.
\end{bridge}

\subsection{Les tabliers des notes de Grothendieck}\label{sec:grdecks}

Trois tabliers reviennent dans les notes de travail de Grothendieck, et les trois arches ci-dessus les lisent ainsi (Lean: \lean{GrothendieckDecks.lean}).
\begin{itemize}
\item La catégorie simpliciale $\Delta$, la première catégorie test de \emph{Pursuing Stacks} (dossiers 108 et 118 du fonds): les ensembles simpliciaux forment un topos à deux valeurs et connexe, puisqu'il existe une application constante entre deux objets quelconques de $\Delta$, et ce n'est pas le topos de préfaisceaux d'un préordre, puisque les deux sommets de $[1]$ sont deux flèches $[0]\to[1]$ (\lean{simplexCategory_twoValued}, \lean{simplexCategory_isConnected}, \lean{simplexCategory_not_spatial}).
\item Un groupe $G$ vu comme groupoïde à un objet, dont le topos de préfaisceaux est $\B G$ (dossier 161-6, p.~22, où Grothendieck ramène une question à ce cas): à deux valeurs, comme le note \cite[IV~8.4.3]{sga4} pour tout monoïde, connexe, et topos de préfaisceaux d'un préordre si et seulement si $G$ est trivial (\lean{singleObj_twoValued}, \lean{singleObj_isConnected}, \lean{singleObj_spatial_iff}).
\item Un ensemble ordonné $I$, dont le topos de préfaisceaux $\widehat I$ est $\mathrm{Top}(I)$ (dossier 37, p.~48; dossier 161-3, p.~26): à deux valeurs si et seulement si deux éléments quelconques sont comparables dans les deux sens, c'est-à-dire si et seulement si son quotient ordonné a au plus un point (\lean{preorder_twoValued_iff}).
\end{itemize}
Ce sont des corollaires immédiats, consignés parce qu'ils sont les premières entrées du registre tirées du fonds.
